\documentclass[letterpaper,10pt]{article}
\usepackage{longtable}
\usepackage{enumerate}
\usepackage{ucs}
\usepackage{multirow}  %%INSTALAR ESTE PAQUETE!!!
\usepackage[latin1]{inputenc}
\usepackage[spanish]{babel}
\usepackage[right=2.2cm,left=2.2cm,top=2.2cm,bottom=2.2cm,headsep=1.2cm,footskip=1.2cm]{geometry}
\usepackage{fancyhdr}
\usepackage{fancybox} 
\usepackage{graphicx,tabularx,latexsym}
\usepackage{amsthm,amsmath,amssymb,amsfonts}
\usepackage{booktabs}
\usepackage{tabulary}
\usepackage{subfigure}
\usepackage{booktabs}
\usepackage{float}
\usepackage{caption}
\usepackage{wrapfig}
\usepackage{listings}
\usepackage{xcolor}
\usepackage[pdftex]{hyperref}
%\usepackage{ragged2e}
\hypersetup
	{
	bookmarks=true, 
	unicode=true,
	pdftitle= Tarea N�4 , 
	pdfauthor="{Aldo Mellado Opazo}", 
	pdfsubject= Programaci�n Orientada al Objeto, 
	pdfnewwindow=true, 
	colorlinks=true, 
	linkcolor=black,
        linktocpage,
            linkcolor=blue,
            anchorcolor=red,
            citecolor=green,
            filecolor=blue,
            menucolor=blue,
            pagecolor=blue,
            urlcolor=blue,
	urlcolor=blue  
	}

\makeindex
\def\figcapskip{1pt}
\def\figtopskip{1pt}
\def\figbottomskip{1pt}
\def\figcapmargin{1pt}
\def\figfarskip{1pt}


\pagestyle{fancy}
 \lhead{Tarea N�4} 
 \rhead{{Programaci�n Orientada al Objeto}} 
 \lfoot{}
 \cfoot{}
 \rfoot{\thepage} 

% ***********************************************************************************
\begin{document}



\renewcommand{\figurename}{\textbf{Fig.}}
\renewcommand{\tablename}{\textbf{Tabla}}
\setlength{\abovecaptionskip}{0pt}
\setlength{\belowcaptionskip}{0pt}


\definecolor{codegreen}{rgb}{0,0.6,0}
\definecolor{codegray}{rgb}{0.5,0.5,0.5}
\definecolor{codepurple}{rgb}{0.58,0,0.82}
\definecolor{backcolour}{rgb}{0.95,0.95,0.92}
 
\lstdefinestyle{mystyle}{
    backgroundcolor=\color{backcolour},   
    commentstyle=\color{codegreen},
    keywordstyle=\color{magenta},
    numberstyle=\tiny\color{codegray},
    stringstyle=\color{codepurple},
    basicstyle=\fontsize{7}{9}\ttfamily,
    breakatwhitespace=false,         
    breaklines=true,                 
    captionpos=b,                    
    keepspaces=true,                 
    numbers=left,                    
    numbersep=5pt,                  
    showspaces=false,                
    showstringspaces=false,
    showtabs=false,                  
    tabsize=2
}

% *************************************************************************************

\begin{titlepage}
\hspace{-1cm} \includegraphics[width=1.5cm]{./imagenes/logo.png}
  \begin{tabular}{l}
  \begin{small} \sc Departamento de Ingenier�a El�ctrica \end{small}\\
 \begin{small} \sc Facultad de Ingenier�a \end{small}\\  
 \begin{small} \sc Universidad de Concepci�n \end{small}\\ 
 \begin{small} \sc Concepci�n, Chile. \end{small}\\                 

  \vspace*{1.5cm}\mbox{}
  \end{tabular}
  %\bigskip

  \vspace*{1.5 in}

\begin{center}
\Huge{{\Huge I}nforme {\Huge N�4}\\}
\vspace{0.3cm}
\large{\bf Simulaci�n de eventos discretos}\\
\vspace{1cm}
{\normalsize \textit{Profesor: Mario Medina}}\\[0.2cm]
{\normalsize \textit{Alumno: Aldo Mellado Opazo}}\\[0.5cm] 
{\normalsize \textit{24 de septiembre de 2018}}
\end{center}

\vspace{9cm}
\begin{center}
\large{Programaci�n Orientada al Objeto}\\
Ingenier�a Civil en Telecomunicaciones
\end{center}
\end{titlepage}


%-----------------------------------------------------------------------------------------------
%----.------------------------------------------------------------------------------------------

\newpage
\setcounter{page}{0}
\thispagestyle{empty}
\tableofcontents


%-----------------------------------------------------------------------------------------------
%----.------------------------------------------------------------------------------------------
%----.------------------------------------------------------------------------------------------
\newpage
\setcounter{page}{1}

\section{Simulaci�n de eventos discretos}

Para esta tarea se entregaba un archivo de texto del cu�l se nos daban tiempos y prioridades, dichos valores deb�an ser procesados y con base en ellos, obtenidos otros par�metros y comportamientos de la red o del sistema que se estaba simulando.\\

Hab�a adem�s como objetivo el simular la respuesta ante los distintos m�todos que se solicitaba.\\

Para la extracci�n de los datos, al comienzo consider� la posibilidad de tomarlos y enviarlos usando mapas o pair a vectores, sin embargo, por simplicidad y por optimizaci�n us� 
\texttt{tuple}, que me permit�a leer desde un archivo de texto, elementos almacenados en forma de columnas y llevarlos a elementos del tipo \texttt{list<float>}, llamados \textit{tiempo1, tiempo2, priority}, respectivamente.\\

Lo comentado se halla en las siguientes l�neas de c�digo.\\

\framebox{
\lstset{style=mystyle, language= C++}
\lstinputlisting[language=C++, firstline=14,lastline=36]{./codes/tarea-4.cpp}
}

\end{document}